\chapter{État de l'art sur la détection du diabète à l'aide des méthodes d'apprentissage automatique}

%----------------------------------------------------------------
\section{Introduction}
%----------------------------------------------------------------

Le diabète sucré (\textit{Diabetes Mellitus}, DM) est une maladie métabolique chronique
caractérisée par une hyperglycémie persistante, c'est-à-dire une concentration anormalement
élevée de glucose dans le sang. Cette anomalie résulte soit d'une déficience de la sécrétion
d'insuline par le pancréas, soit d'une résistance des cellules à l'action de cette hormone,
soit d'une combinaison des deux mécanismes~\citep{iparraguirre2023}. L'insuline, hormone
produite par les cellules bêta des îlots de Langerhans, joue un rôle fondamental dans la
régulation glycémique : elle permet aux cellules d'absorber le glucose circulant dans le sang
pour le convertir en énergie. En son absence ou en cas de résistance à son action, le glucose
s'accumule dans la circulation sanguine, entraînant à long terme des dommages irréversibles
sur de nombreux organes et systèmes.

\subsection*{Types de diabète}

On distingue principalement trois formes cliniques du diabète :

\textbf{Le diabète de type 1 (DT1)} est une maladie auto-immune dans laquelle le système
immunitaire attaque et détruit les cellules bêta pancréatiques productrices d'insuline. Il
représente environ 5 à 10\% des cas diagnostiqués et se manifeste généralement durant
l'enfance ou l'adolescence. Les patients atteints de DT1 dépendent à vie d'injections
d'insuline exogène pour leur survie~\citep{iparraguirre2023, zou2018}.

\textbf{Le diabète de type 2 (DT2)} constitue la forme la plus répandue, représentant plus
de 90\% des cas mondiaux. Il se caractérise par une résistance progressive à l'insuline,
combinée à une sécrétion pancréatique insuffisante pour compenser cette résistance. Son
développement est fortement lié au mode de vie — sédentarité, obésité, alimentation
hypercalorique — ainsi qu'à une prédisposition génétique. Contrairement au DT1, il survient
le plus souvent à l'âge adulte, bien qu'une tendance croissante soit observée chez les
adolescents dans les pays développés~\citep{zou2018}.

\textbf{Le diabète gestationnel} apparaît durant la grossesse chez des femmes ne présentant
pas de diabète préalable. Il se résorbe généralement après l'accouchement, mais augmente
significativement le risque de développer un DT2 ultérieurement, aussi bien chez la mère
que chez l'enfant. La base de données Pima Indians Diabetes Database (PIDD), utilisée dans
ce travail, concerne exclusivement des femmes adultes de la communauté Pima, une population
présentant une prévalence exceptionnellement élevée de DT2, dont une part importante est
liée à des antécédents de diabète gestationnel.

\subsection*{Facteurs de risque et variables cliniques déterminantes}

La progression vers le diabète de type 2 est influencée par un ensemble de facteurs
biologiques et comportementaux. Parmi les principaux indicateurs cliniques figurant dans la
littérature et dans la PIDD, on trouve~\citep{zou2018, iparraguirre2023} :

\begin{itemize}
	\item \textbf{La glycémie (glucose plasmatique)} : principal marqueur diagnostique, mesurée
	à jeun ou après une charge orale en glucose. Un taux supérieur à 126 mg/dL à jeun est
	diagnostique du diabète selon les critères de l'American Diabetes Association (ADA). C'est
	la variable de plus grande importance prédictive dans les études sur la PIDD.
	
	\item \textbf{L'indice de masse corporelle (IMC ou BMI)} : fortement corrélé au risque de DT2,
	un IMC supérieur à 30 kg/m\textsuperscript{2} indique une obésité et constitue un facteur
	de risque majeur de résistance à l'insuline.
	
	\item \textbf{Le taux d'insuline} : une élévation chronique de l'insulinémie peut signaler
	une résistance à l'insuline précédant le diabète clinique.
	
	\item \textbf{La pression artérielle diastolique} : l'hypertension est à la fois un facteur
	de risque et une comorbidité fréquente du diabète, augmentant le risque cardiovasculaire.
	
	\item \textbf{L'épaisseur du pli cutané tricipital} : indicateur indirect de la masse grasse
	corporelle, corrélé à l'obésité et à la résistance à l'insuline.
	
	\item \textbf{La fonction pedigree du diabète (Diabetes Pedigree Function)} : score composite
	reflétant l'hérédité familiale diabétique, intégrant la prévalence du diabète chez les
	apparentés et leur degré de parenté.
	
	\item \textbf{L'âge} : le risque de DT2 augmente avec l'âge, particulièrement après 45 ans.
	
	\item \textbf{Le nombre de grossesses} : dans la population Pima, le nombre de grossesses est
	corrélé à l'âge et constitue un facteur de risque indirect via le diabète gestationnel.
\end{itemize}

\subsection*{Conséquences et enjeux cliniques}

Non traité ou mal contrôlé, le diabète entraîne des complications graves : néphropathie
diabétique pouvant évoluer vers l'insuffisance rénale terminale, rétinopathie et cécité,
neuropathie périphérique, maladies cardiovasculaires et artériopathie des membres inférieurs
pouvant mener à l'amputation. La détection précoce est donc cruciale : elle permet d'initier
des interventions sur le mode de vie et des traitements pharmacologiques avant l'apparition
des complications irréversibles~\citep{wee2024, khanam2021}.

%----------------------------------------------------------------
\section{Approches traditionnelles de détection du diabète}
%----------------------------------------------------------------

Historiquement, le diagnostic du diabète repose sur des critères biochimiques standardisés
établis par des organisations de référence telles que l'OMS et l'ADA. Ces critères incluent
la mesure de la glycémie à jeun ($\geq$ 126 mg/dL), la glycémie deux heures après une charge
orale en glucose de 75g ($\geq$ 200 mg/dL), ou le taux d'hémoglobine glyquée HbA1c
($\geq$ 6,5\%). Ces mesures constituent des seuils fixes, appliqués uniformément à l'ensemble
de la population.

Si ces méthodes présentent l'avantage d'être standardisées, reproductibles et bien comprises
cliniquement, elles souffrent de plusieurs limitations importantes dans un contexte de
prévention et de dépistage de masse~\citep{sisodia2018} :

\begin{enumerate}
	\item \textbf{Absence d'intégration multi-variables} : les critères traditionnels ne tiennent
	pas compte simultanément des multiples facteurs de risque. Or, la progression vers le DT2
	est un phénomène multidimensionnel impliquant des interactions complexes entre variables
	génétiques, métaboliques et comportementales.
	
	\item \textbf{Détection tardive} : les critères actuels diagnostiquent un diabète déclaré,
	souvent après plusieurs années de dysrégulation glycémique silencieuse (stade prédiabétique).
	
	\item \textbf{Coût et accessibilité} : les tests biologiques (HOMA-IR, OGTT) requièrent des
	prélèvements à jeun et une infrastructure de laboratoire, ce qui limite leur accessibilité
	dans les systèmes de santé à ressources limitées.
	
	\item \textbf{Variabilité inter-individuelle} : deux individus présentant la même glycémie à
	jeun peuvent avoir des profils de risque très différents selon leurs caractéristiques
	anthropométriques, héréditaires et comportementales.
\end{enumerate}

Ces limites ont motivé l'exploration de méthodes analytiques capables d'exploiter l'ensemble
des données cliniques disponibles pour produire des prédictions personnalisées du risque
diabétique~\citep{sisodia2018, kavakiotis2017}.

%----------------------------------------------------------------
\section{Méthodes basées sur l'apprentissage automatique}
%----------------------------------------------------------------

L'application de l'apprentissage automatique (\textit{machine learning}, ML) à la détection
du diabète a connu une croissance exponentielle depuis les années 2010. Ces méthodes
présentent plusieurs avantages structurels : capacité à modéliser des relations non linéaires
entre variables, aptitude à traiter des jeux de données à forte dimensionnalité, et possibilité
d'intégrer des caractéristiques hétérogènes dans un cadre unifié de décision~\citep{wee2024}.

\subsection{Régression logistique et méthodes linéaires}

La régression logistique (RL) constitue souvent le modèle de référence (\textit{baseline}) dans
les études de classification médicale. Pour la prédiction du diabète, la RL modélise la
probabilité d'appartenance à la classe ``diabétique'' en fonction d'une combinaison linéaire
des variables d'entrée, transformée par la fonction sigmoïde~\citep{iparraguirre2023}. Bien
que simple et interprétable, ce modèle suppose une séparabilité linéaire entre les classes,
hypothèse rarement vérifiée dans les données cliniques complexes. Les performances rapportées
sur la PIDD se situent généralement entre 72\% et 78\% d'exactitude selon les études.

\subsection{Arbres de décision et méthodes ensemblistes basiques}

Les arbres de décision (Decision Trees, DT) ont été parmi les premiers algorithmes appliqués
à la détection du diabète sur la PIDD~\citep{sisodia2018}. Leur interprétabilité naturelle
— chaque décision pouvant être visualisée sous forme d'un arbre binaire — les rend
particulièrement attractifs pour les applications médicales. Cependant, leur forte sensibilité
au surapprentissage limite leurs performances en généralisation. \citeauthor{zou2018}
montrent que dans un arbre de décision entraîné sur la PIDD, le glucose constitue le nœud
racine, confirmant ainsi la pertinence clinique de ce prédicteur~\citep{zou2018}.

\subsection{Méthodes à noyaux : les SVM}

Les Machines à Vecteurs de Support (SVM) ont été appliquées à la détection du diabète dès
les travaux pionniers de la fin des années 1990. Elles présentent l'avantage de trouver un
hyperplan de séparation optimal dans des espaces de grande dimension, avec une bonne
résistance au surapprentissage grâce à la maximisation de la marge. Sur la PIDD, les SVM
avec noyau RBF (\textit{Radial Basis Function}) obtiennent généralement des précisions entre
71\% et 82\% selon les stratégies de prétraitement adoptées~\citep{iparraguirre2023,
	wee2024}. Leur sensibilité à l'échelle des données et au choix de l'hyperparamètre $C$
nécessite une normalisation minutieuse et une optimisation des hyperparamètres.

\subsection{Méthodes probabilistes et apprentissage bayésien}

Les classifieurs Naïf de Bayes (\textit{Naïve Bayes}, NB) et ses variantes (Bernoulli NB,
Gaussian NB) offrent une approche probabiliste basée sur le théorème de Bayes, sous
l'hypothèse d'indépendance conditionnelle entre les variables. Malgré la relative simplic\-ité
de cette hypothèse, ces méthodes produisent des résultats compétitifs sur des jeux de données
de taille modérée~\citep{iparraguirre2023}. \citeauthor{iparraguirre2023} rapportent une
exactitude de 77,2\% avec le modèle Bernoulli NB couplé à la technique SMOTE pour
rééquilibrer les classes.

\subsection{Méthodes ensemblistes avancées : Random Forest et XGBoost}

Les méthodes ensemblistes représentent l'état de l'art pour la classification sur données
tabulaires. \citeauthor{zou2018} démontrent dans une étude comparative pionnière que les
forêts aléatoires atteignent la meilleure exactitude (ACC = 80,84\%) sur la PIDD lorsque
l'ensemble des attributs est utilisé, surpassant les réseaux de neurones et les arbres de
décision isolés~\citep{zou2018}. Ces résultats ont été confirmés et approfondis par de
nombreux travaux ultérieurs~\citep{wee2024, gundogdu2023, khanam2021}.

XGBoost a progressivement émergé comme l'algorithme le plus compétitif dans ce domaine.
\citeauthor{gundogdu2023} proposent une architecture combinant Random Forest pour la
sélection de variables et XGBoost pour la classification, atteignant une AUC de 98,6\% sur
un jeu de données de 520 patients diabétiques du Bangladesh~\citep{gundogdu2023}. Ce
résultat illustre l'intérêt de combiner les forces des deux méthodes : la robustesse de la
Random Forest pour l'identification des variables pertinentes, et la précision prédictive du
gradient boosting.

%----------------------------------------------------------------
\section{Études comparatives utilisant SVM, Random Forest et XGBoost}
%----------------------------------------------------------------

Un corpus croissant de travaux propose des comparaisons directes entre les trois algorithmes
ciblés dans ce mémoire, offrant une base empirique solide pour orienter le choix du modèle
en contexte clinique.

\subsection{Comparaisons sur la base PIDD}

\citeauthor{iparraguirre2023} comparent cinq algorithmes — K-NN, Naïf de Bayes, arbres de
décision, régression logistique et SVM — sur la PIDD avec application de SMOTE pour
l'équilibrage des classes~\citep{iparraguirre2023}. Leurs résultats montrent que K-NN et
Bernoulli NB surpassent le SVM (71,7\% d'exactitude) dans ce contexte. Cette relative
contre-performance du SVM est attribuée à la petite taille du jeu de données et à la
sensibilité de l'algorithme aux hyperparamètres dans des espaces de faible dimensionnalité.

\citeauthor{zou2018} proposent une comparaison systématique sur la PIDD incluant les
forêts aléatoires, les arbres de décision et les réseaux de neurones, avec application de
techniques de réduction de dimensionnalité (PCA et mRMR)~\citep{zou2018}. La Random
Forest obtient la meilleure exactitude globale (80,84\%), ce résultat étant robuste aux
différentes configurations de sélection de variables testées. Les auteurs soulignent que le
glucose est systématiquement identifié comme la variable de plus haute importance, tant par
les critères d'entropie dans les arbres que par l'importance des variables dans la Random
Forest.

\subsection{Comparaisons multi-algorithmes incluant XGBoost}

\citeauthor{wee2024} proposent une revue approfondie des méthodes de ML et de deep
learning appliquées à la détection du diabète sur plusieurs bases de données publiques~\citep{wee2024}. Leur analyse met en évidence la supériorité systématique des méthodes
ensemblistes (Random Forest, XGBoost) sur les méthodes à modèle unique (régression
logistique, SVM) en termes d'exactitude, de F1-score et d'AUC-ROC. XGBoost se distingue
notamment par ses mécanismes intégrés de régularisation et sa gestion native des valeurs
manquantes, deux propriétés particulièrement avantageuses pour les données médicales
réelles souvent incomplètes.

\citeauthor{khanam2021} réalisent une comparaison de sept algorithmes de ML pour la
prédiction du diabète, incluant SVM, Random Forest, régression logistique, K-NN, arbres de
décision, réseaux de neurones et Naïf de Bayes~\citep{khanam2021}. Sur la PIDD, les résultats
montrent que la régression logistique et le SVM obtiennent les meilleures performances en
termes d'exactitude, tandis que la Random Forest excelle sur le rappel (\textit{recall}),
métrique particulièrement importante en contexte médical pour minimiser les faux négatifs
(patients diabétiques non détectés).

\subsection{Apport des techniques de prétraitement et d'optimisation}

Plusieurs études soulignent que les performances comparatives dépendent fortement des
stratégies de prétraitement adoptées. \citeauthor{gundogdu2023} montrent que la sélection
de variables par Random Forest permet d'améliorer significativement les performances de
XGBoost en réduisant le bruit introduit par les variables non informatives~\citep{gundogdu2023}.
\citeauthor{iparraguirre2023} démontrent que l'application de SMOTE améliore systématiquement
les métriques de classification pour tous les algorithmes testés sur la PIDD, en réduisant
le déséquilibre entre les classes positives (diabétiques, 35\% du jeu de données) et
négatives~\citep{iparraguirre2023}.

Le Tableau~\ref{tab:comparaison_etudes} synthétise les principales études comparatives
identifiées dans la littérature.

\begin{table}[htbp]
	\centering
	\caption{Synthèse des principales études comparatives de ML pour la détection du diabète}
	\label{tab:comparaison_etudes}
	\small
	\begin{tabular}{p{2.8cm} p{2.8cm} p{2.5cm} p{3cm} p{2cm}}
		\hline
		\textbf{Étude} & \textbf{Dataset} & \textbf{Algorithmes} & \textbf{Meilleur résultat} & \textbf{Métrique} \\
		\hline
		Zou et al. \cite{zou2018} & PIDD + Luzhou & RF, DT, NN & RF: 80,84\% & ACC \\
		\hline
		Iparraguirre et al. \cite{iparraguirre2023} & PIDD (SMOTE) & K-NN, NB, DT, LR, SVM & K-NN: 79,6\% & ACC \\
		\hline
		Wee et al. \cite{wee2024} & Multiple & RF, XGB, SVM, DL & XGBoost & AUC-ROC \\
		\hline
		Khanam \& Foo \cite{khanam2021} & PIDD & 7 algorithmes & LR/SVM & ACC / Recall \\
		\hline
		Gündoğdu \cite{gundogdu2023} & Sylhet, BD & RF + XGB & RF-XGB: 98,6\% & AUC \\
		\hline
		Sisodia \& Sisodia \cite{sisodia2018} & PIDD & NB, DT, SVM & NB: 76,3\% & ACC \\
		\hline
	\end{tabular}
\end{table}

%----------------------------------------------------------------
\section{Analyse critique des travaux existants}
%----------------------------------------------------------------

L'examen de la littérature permet d'identifier plusieurs constats et limites récurrents qui
méritent une discussion approfondie.

\subsection{Hétérogénéité des protocoles expérimentaux}

La comparabilité des résultats entre études est sérieusement compromise par l'hétérogénéité
des protocoles. Les différences portent sur la stratégie de partitionnement des données
(hold-out simple vs. validation croisée k-fold), le traitement des valeurs manquantes et des
zéros biologiquement impossibles dans la PIDD (certaines études les suppriment, d'autres les
imputent), l'usage ou non de techniques de rééchantillonnage (SMOTE, ADASYN), ainsi que le
choix des métriques d'évaluation. Ces disparités rendent les comparaisons directes entre
études délicates et expliquent en partie la grande variabilité des performances rapportées
pour un même algorithme~\citep{wee2024}.

\subsection{Déséquilibre des classes et biais des métriques}

La PIDD présente un déséquilibre significatif entre les classes : environ 35\% de cas
positifs contre 65\% de cas négatifs. Cette asymétrie induit un biais dans l'interprétation
de l'exactitude simple (\textit{accuracy}), qui peut être artificiellement élevée même pour
un classifieur majoritaire. Plusieurs études publiées antérieurement ne reportent que
l'exactitude, occultant les performances sur la classe minoritaire (diabétiques) qui est
pourtant la plus cliniquement pertinente~\citep{khanam2021, iparraguirre2023}. Une analyse
rigoureuse doit systématiquement inclure le rappel (sensibilité), la précision, le F1-score
et l'AUC-ROC pour chaque classe.

\subsection{Généralisation limitée}

La majorité des études comparatives se concentrent sur la PIDD, une base de données de
768 instances issues d'une population très spécifique (femmes Pima de plus de 21 ans).
La généralisation de ces résultats à d'autres populations reste une question ouverte. Les
caractéristiques démographiques, génétiques et métaboliques de la population Pima sont
atypiques à l'échelle mondiale, ce qui limite l'extrapolation directe des modèles
développés~\citep{zou2018, wee2024}.

\subsection{Explicabilité insuffisante}

Peu d'études analysent l'importance des variables de manière rigoureuse ou n'explorent les
mécanismes de décision des modèles au-delà des simples scores d'importance. L'émergence
d'outils d'explicabilité comme SHAP (\textit{SHapley Additive exPlanations}) offre des
perspectives prometteuses pour rendre les modèles ensemblistes interprétables par les
cliniciens, mais leur utilisation reste marginale dans la littérature sur la PIDD~\citep{wee2024}.

\subsection{Optimisation du seuil de décision}

La quasi-totalité des études utilise le seuil de classification par défaut de 0,5, sans
l'optimiser en fonction du contexte clinique. En détection du diabète, un faux négatif
(patient diabétique non détecté) est généralement plus coûteux qu'un faux positif, ce qui
justifie l'utilisation de seuils inférieurs à 0,5 pour maximiser la sensibilité. Cette
dimension de l'évaluation est systématiquement négligée dans les études
existantes~\citep{khanam2021}.

%----------------------------------------------------------------
\section{Positionnement et apport du présent travail}
%----------------------------------------------------------------

À la lumière de cette revue de littérature, le présent travail se positionne dans le
prolongement des études comparatives existantes tout en apportant plusieurs contributions
spécifiques.

\textbf{Premièrement}, contrairement aux études qui comparent un grand nombre d'algorithmes
superficiellement, ce mémoire se concentre sur trois algorithmes majeurs — SVM, Random
Forest et XGBoost — avec une analyse approfondie de leurs mécanismes théoriques et de leur
comportement face aux spécificités de la PIDD. Ce choix est motivé par la complémentarité
des approches : les SVM représentent les méthodes à noyau, la Random Forest représente le
bagging ensembliste, et XGBoost représente le boosting par gradient.

\textbf{Deuxièmement}, l'évaluation adoptée dans ce travail va au-delà de la simple
exactitude pour inclure un ensemble complet de métriques adapté au contexte médical :
précision, rappel, F1-score, AUC-ROC et matrice de confusion par classe. Une attention
particulière est portée à la sensibilité (rappel de la classe positive), directement liée
au taux de détection des patients diabétiques.

\textbf{Troisièmement}, ce mémoire intègre une analyse de l'importance des variables pour
chaque modèle, permettant d'identifier les prédicteurs les plus déterminants dans la PIDD
et de confronter les conclusions algorithmiques aux connaissances médicales établies.

\textbf{Quatrièmement}, la gestion du déséquilibre des classes est traitée de manière
explicite, en comparant différentes stratégies : pondération des classes, techniques de
rééchantillonnage, et paramètres natifs (comme \texttt{scale\_pos\_weight} dans XGBoost)
pour évaluer leur impact sur les performances finales.

En somme, ce travail vise à produire une évaluation comparative rigoureuse, reproductible
et cliniquement orientée, dont les conclusions peuvent guider des praticiens ou des équipes
de recherche souhaitant déployer un modèle de détection du diabète sur la PIDD ou des
jeux de données cliniques similaires.

